\section{Conclusiones}
\label{sec:conclusiones}

Siguiendo la guía de OpenHPC 4.1 para Rocky Linux 10 con Warewulf y Slurm, logramos instalar el nodo maestro y preparar los nodos de cómputo sobre máquinas virtuales. Al nodo maestro se le instaló el sistema operativo base con SELinux deshabilitado, el servidor NTP, los componentes de OpenHPC, el servidor de Slurm y Warewulf. A partir de él se generó la imagen \texttt{rocky-10} que reciben los nodos de cómputo.

El objetivo principal se cumplió en parte. Los nodos de cómputo llegaron a arrancar por red con la imagen distribuida por Warewulf y recibieron la dirección que les correspondía en la red privada, como muestran la \cref{fig:nodo-login} y la \cref{fig:maestro-dhcp}. Sin embargo, no alcanzamos a ejecutar a tiempo la prueba del hola mundo con MPI, por lo que no pudimos comprobar que el cluster ejecute trabajos en paralelo a través de Slurm.

La mayor dificultad estuvo en la red y no en la instalación de los paquetes. Durante varios intentos el servidor DHCP del nodo maestro ofrecía la dirección a cada nodo, pero estos la rechazaban y no arrancaban la imagen por red, incluso con la dirección asignada manualmente. De esto concluimos que en un cluster y la configuración de las interfaces, del direccionamiento de la red privada y del servicio DHCP debe quedar perfectamente configurado ya que si no atrasa bastante la instalación del Cluster debido a que sin red no hay comunicación entre los nodos ni imagen para los mismos.

Trabajar con máquinas virtuales nos permitió repetir la instalación y probar varios programas de virtualización (VirtualBox, KVM y VMware), aunque cada uno maneja de forma distinta los adaptadores de red, y eso influye directamente en el arranque por red de los nodos, gracias a esto terminamos optando por KVM ya que tiene más opciones de configuración de red.

\subsection{Conclusiones personales}
Mauricio:
El haber realizado la instalación de Rocky y del nodo maestro en VirtualBox más aparte haber tenido que realizar llaves y firmarlas debido a problemas entre el kernel y el mismo software nos atrase mucho más de lo que debimos de haber tardado, finalmente ya que cambiamos a software que era VMware fue más fácil la instalación y como ya conocía el proceso de firmar las llaves del kernel fue más ágil, pero aún así no conseguí que funcionará ni como nodo maestro ni como nodo de cómputo con el nodo maestro que ya funcionaba de KVM, y fue por ello que mejor entre todo el equipo nos movimos a KVM ya que en fedora es mucho más configurable y fácil de manejar que todos los softwares anteriores, siendo este último el único que funciono y sin muchas trabas fuera del mismo problema del DHCP, considero que hubieramos terminado a tiempo si hubieramos tenido softwares ya definidos como para no hacer prueba y error.
\\[1em]
Jacobo:
Tres de los cuatro integrantes del equipo contamos con el sistema operativo de Fedora, por lo que al inicio tuvimos algunas dificultades con el software de Virtualbox por lo que finalmente decidimos usar QEMU/KVM en su lugar ya que al parecer este software esta mejor optimizado para este S.O. en especifico, sin embargo esto también nos trajo algunos problemas, el principal fue la conexion Ethernet ya que era algo mas difícil usar adaptadores tipo pueden en las VM, sin embargo al final pudimos resolver este problem usando una NIC MacVTap.
Finalmente puedo decir que esta practica fue posiblemente una de las mas difíciles a las que me he enfrentado dado que es necesario tener conocimientos en Linux con sistemas Red Hat como lo es Rocky, ademas de conocimientos básicos en redes, pero no lo normal si no precisamente en temas básicos como lo son los tipos de interfaces: NAT, puente, etc y finalmente conocimientos básicos sobre maquinas virtuales. Sin duda este experimento me trajo muchos conocimientos que estoy seguro aplicare en proyectos futuros.
\\[1em]
Francisco:
La realización de este proyecto me permitió comprender mejor el proceso de configuración e implementación de un clúster, así como las dificultades que pueden presentarse durante su desarrollo. A lo largo del proyecto surgieron diversos problemas relacionados con las máquinas virtuales, la configuración de las interfaces de red, la instalación de herramientas y la comunicación entre los nodos. A pesar de los intentos realizados y de las diferentes soluciones que se investigaron, no fue posible establecer la conexión final del clúster. Sin embargo, considero que esta experiencia fue importante para mi aprendizaje, ya que me permitió desarrollar habilidades para identificar errores, investigar posibles soluciones y comprender la importancia de realizar cada configuración correctamente. Personalmente, aprendí que este tipo de proyectos requiere paciencia, tiempo, organización y trabajo en equipo. Aunque no se logró cumplir completamente el objetivo, los conocimientos adquiridos durante el proceso representan una base útil para futuros proyectos, en los que podré aplicar lo aprendido y enfrentar problemas similares con mayor experiencia y preparación.
\\[1em]
Guillermo:
La implementación del clúster con OpenHPC y Rocky Linux fue un desafío complejo pero muy interesante. Quedó claro que una mejor planeación inicial habría agilizado considerablemente nuestro avance. Interconectar varias máquinas virtuales es complicado, y lidiar con los distintos sistemas y configuraciones en la PC de cada integrante lo dificultó más. Lo ideal sería estandarizar el entorno para todos, partiendo de una instalación de sistema operativo desde cero.
En mi caso, instalaciones previas de Docker requirieron ajustes extra para lograr tener red en la máquina virtual. A esto se suma que el manual oficial es excesivamente técnico y asume muchos pasos y comandos de configuración. Finalmente, tras la experiencia, creo que existen distribuciones basadas en Red Hat más adecuadas que Rocky Linux para este tipo de implementaciones.
